\section{Epistemic and Other Reasoning Modes}\label{sec:epistemic}

The substrate generalizes beyond what is \emph{true} or \emph{probable}. This section covers three further modes that reuse the same CUDA join, circuit, and CDCL machinery---epistemic reasoning over world views, probabilistic aggregates, and well-founded semantics---broadening what a neural network can be integrated with while keeping data-plane execution GPU-backed.

\subsection{Modal Surface and the Epistemic IR}

Epistemic logic programming asks what is \emph{known} or \emph{possible} across the multiple stable models a program admits---the natural setting for introspection and reasoning under incomplete knowledge. A program may use four modal literals over a predicate---\texttt{know}, \texttt{possible}, \texttt{not know}, \texttt{not possible}. Rather than rewrite them away at parse time, the compiler preserves them in an Epistemic IR (EIR) between the AST and execution planning. The high-level dispatcher classifies modal dependencies before selecting Generate--Propagate--Test, a mode-specific positive-cycle fixpoint, or GPU-backed WFS. The accepted surface spans finite nested modal chains, epistemic integrity constraints (which prune candidate world views on the GPU), rules that couple several epistemic predicates, same-name multi-arity modal predicates, and recursive epistemic execution.

\subsection{World-View Semantics}

A world view is a non-empty set of stable models. \texttt{know p} holds when \texttt{p} appears in every model of the world view; \texttt{possible p} holds when it appears in at least one. Two semantics are selectable. Under the default FAEEL (Founded Autoepistemic Equilibrium Logic), a circular derivation contributes a tuple only when independent founded support exists. A program containing only \texttt{p :- possible p.} therefore executes to an empty founded extension rather than failing with an unsupported-construct diagnostic. A Gelfond-1991-style compatibility mode admits that self-supporting tuple. Non-epistemic programs evaluate identically under both.

\subsection{Epistemic Execution Routes}\label{subsec:epistemic-gpu}

After EIR construction, acyclic modal programs follow a Generate--Propagate--Test schedule whose phases run as CUDA kernels. \emph{Generate} enumerates bounded candidate assumptions as device bitsets; \emph{propagate} prunes candidates that contradict known or rejected facts; \emph{test} validates survivors against the program's world views, checking stable-model tuple membership on-device per arity. Validating a candidate on this route reuses the on-GPU CDCL solver already built for compilation verification (\Cref{sec:prob}), the same substrate that serves satisfiability, bounded MaxSAT and assumption-based queries, and accepted world-view evidence feeds the existing GPU-native exact path.

A positive FAEEL dependency cycle reduces to ordinary rules and reaches the founded least fixpoint in the existing recursive engine. A supported Gelfond-1991 cycle takes a different plan: its selected positive \texttt{possible} gates must use the same terms as their rule heads and remain within one recursive dependency component. The runtime first relaxes those gates to compute an upper bound, then reevaluates from the extensional inputs while each gate reads a frozen snapshot from the preceding pass, and GPU set comparison over all intensional relations detects the descending greatest compatibility fixpoint. Consequently, predicate-level recursion cannot preserve a tuple unless the recursive rules return that concrete tuple, and candidate relations with disjoint tuple domains descend to the empty extension.

Supported cycles through negation use the GPU-backed WFS alternating fixpoint. An unrelated WFS component can execute inside a Gelfond-1991 pass, but recursive negation or aggregation within the compatibility component is rejected, because either would make the descending operator nonmonotone.

The reduced ordinary work reuses the optimizer, WCOJ planner and helper-splitting machinery of \Cref{sec:datalog}, and \emph{epistemic splitting} decomposes eligible acyclic programs into independent components solved separately and recombined. No CPU fallback is provided for the accepted data-plane hot paths: a construct these routes do not support is rejected with a typed error, not served by a host implementation. Transfer budgets track data movement; the host sequences Gelfond-1991 refinement and WFS iterations, while relation evaluation, snapshot copies and convergence comparisons remain GPU-backed.

\subsection{Probabilistic Aggregates}

Finite aggregate queries (\texttt{count}, \texttt{sum}, \texttt{min}, \texttt{max}, \texttt{logsumexp}) are supported in both exact provenance/PIR inference and Monte Carlo execution, subject to a typed exact-domain cap that rejects domains too large for tractable enumeration. For small finite \texttt{count} domains, \emph{aggregate lifting} replaces naive world enumeration with exact cardinality dynamic programming over a compact state set---collapsing, for example, a $2^{17}$-outcome enumeration into a few hundred lifted states with identical results---and accepted \texttt{count} aggregates dispatch a GPU-native exact evaluator rather than a CPU finite-world shortcut, keeping the path device-resident.

\subsection{Well-Founded Semantics}

Programs that violate stratification---such as the classic \texttt{p :- not q.\ q :- not p.}---do not have a unique two-valued least-fixpoint interpretation. For supported shapes, xlog instead assigns three-valued truth (true, false, undefined) through Well-Founded Semantics, computing an alternating fixed point: mark the greatest unfounded set false, derive consequences true, and repeat until stable. For probabilistic programs, true atoms receive normal probability and gradient while false and undefined atoms are assigned probability zero, matching the conservative treatment of non-stratified programs. Probabilistic provenance analysis and epistemic dependency dispatch select the GPU-backed WFS route for the nonmonotone components they support; other cyclic shapes return a typed diagnostic.
